\section{Distance and the Emergence of Space}
\label{sec:space}

\subsection{Decomposition of $C_{ij}$}

The causal log matrix is generically asymmetric.
We decompose it into symmetric and antisymmetric parts:

\begin{equation}
  S_{ij} = \tfrac{1}{2}(C_{ij} + C_{ji}),
  \qquad
  A_{ij} = \tfrac{1}{2}(C_{ij} - C_{ji}).
  \label{eq:decomp}
\end{equation}

$S_{ij}$ encodes mutual connectivity (the spatial component);
$A_{ij}$ encodes directed flow (the temporal and causal component).
Time and space are not independently postulated;
they arise as two aspects of the same object $C_{ij}$.

\subsection{Spatial Distance}

For a path $\gamma: i \to j$ through the network,
define its \textbf{propagation strength}:

\[
  \tilde{S}_{ij}
  = \max_{\gamma: i \to j} \prod_{(ab)\in\gamma} S_{ab}.
\]

The \textbf{distance} between nodes is:

\begin{keybox}
\begin{equation}
  d_{ij} = -\log\,\tilde{S}_{ij}
  = \min_{\gamma: i \to j} \sum_{(ab)\in\gamma} (-\log S_{ab}).
  \label{eq:distance}
\end{equation}
\end{keybox}

\noindent
This is a shortest-path problem on a weighted network.
The logarithm converts products to sums and ensures that distance accumulates
additively along paths.

\begin{proposition}
$d_{ij}$ satisfies non-negativity and the triangle inequality,
inherited from the path structure of the network.
\end{proposition}

\begin{keybox}
\textit{Distance is resistance to influence propagation.
Space is not a background arena; it is a measure of how difficult
information finds it to propagate between nodes.}
\end{keybox}

\subsection{Coarse-Graining and the Metric Tensor}

In a coarse-grained limit, grouping nodes into spatial blocks and averaging
produces a smooth effective distance:

\[
  d(x, y)^2 \approx g_{\mu\nu}(x)\,(x^\mu - y^\mu)(x^\nu - y^\nu).
\]

\begin{keybox}
\begin{equation}
  g_{\mu\nu} = \text{effective description of } d_{ij}
  \text{ under coarse-graining.}
  \label{eq:metric}
\end{equation}
\end{keybox}

\noindent
The metric tensor $g_{\mu\nu}$ is not a fundamental field.
It is the coarse-grained approximation of the log-distance structure.
No coordinate system is assumed; coordinates label coarse-grained blocks.

\subsection{Time and Space Unified}

The two emergent structures can be summarized as:

\begin{keybox}
\begin{align}
  \frac{d\tau_i}{dt} &= \rho_i
  \quad &&\text{(time: quantity of logs)}
  \label{eq:time_summary}\\[4pt]
  d_{ij} &= -\log\,\tilde{S}_{ij}
  \quad &&\text{(space: structure of logs)}
  \label{eq:space_summary}
\end{align}
\end{keybox}

\noindent
Time and space originate from different aspects of the same
underlying object $C_{ij}$.
